% W5 outcome: exact higher pair--Palm reduction, sharp deterministic bounds,
% and a complete audit of the proposed double-Cauchy--Schwarz route.
% No unconditional proof of (HM)_7 or (HM)_8 is claimed.

\subsection*{Higher pair--Palm moments: exact reduction and obstruction}

\paragraph{Two normalization corrections.}
The high-moment hypothesis in
Section~\ref{ssec:high-moment}, and the proved estimate
Proposition~\ref{prop:hm2}, use
\begin{equation}\label{eq:pairpalm-lambda}
 \lambda=\lambda_X
 =\sum_{p\in\mathcal P_X}\frac{Z(p)}p,
 \qquad
 \mathcal P_X=\{p\text{ prime}:X<p\leq2X\}.
\end{equation}
Thus the unweighted expression $\sum_p1/p$ in the task description
must be replaced by~\eqref{eq:pairpalm-lambda}; otherwise even the
scale of Proposition~\ref{prop:hm2} is not the one stated in the paper.
There is a second, independent multiplicity issue.  The third-prime
quantity in the pair--Palm description must count \emph{all} third
primes which hit, not merely indicate that one exists.  At a row with
$K_X(m)=4$, the existential version contributes $(4)_2=12$, whereas
the required ordered third factorial moment is $(4)_3=24$.

Put $M=X^2$ and, as in Theorem~\ref{thm:hm3-palm-characterization}, set
\[
 \Omega_p=\{0\leq m<M:m\bmod p\in\mathcal Z_p\},
 \qquad f_p(m)=\mathbf1_{\Omega_p}(m),
 \qquad K(m)=\sum_{p\in\mathcal P_X}f_p(m).
\]
For $k\geq2$ write
\[
 F_k(X)=\sum_{m<M}(K(m))_k.
\]

\begin{lemma}[A four-constant second-moment bound]
\label{lem:pairpalm-hm2-four}
One has
\begin{equation}\label{eq:pairpalm-hm2-four}
 F_2(X)\leq4M\lambda^2.
\end{equation}
\end{lemma}

\begin{proof}
For distinct $p,q\in\mathcal P_X$ one has $pq>X^2=M$.
Consequently a fixed residue pair modulo $(p,q)$ has at most one
representative in $[0,M)$, and hence
\[
 |\Omega_p\cap\Omega_q|\leq Z(p)Z(q).
\]
Summing over ordered pairs and using $p\leq2X$ gives
\[
 F_2(X)
 \leq\left(\sum_pZ(p)\right)^2
 \leq\left(2X\sum_p\frac{Z(p)}p\right)^2
 =4M\lambda^2.
\]
This only sharpens the harmless constant~$5$ in
Proposition~\ref{prop:hm2}; no new distributional input is used.
\end{proof}

For ordered $p\ne q$ define
\[
 J_{p,q}=|\Omega_p\cap\Omega_q|,
 \qquad
 \lambda_{p,q}
 =\sum_{\ell\in\mathcal P_X\setminus\{p,q\}}
       \frac{Z(\ell)}\ell.
\]
For a fixed $k\geq3$, put $r=k-2$ and define the ordered
$r$-extension count
\begin{align}
 T_{p,q}^{[r]}
 &=\sum_{m\in\Omega_p\cap\Omega_q}(K(m)-2)_r
                                                        \label{eq:pairpalm-T}\\
 &=\sum_{\substack{
     (\ell_1,\ldots,\ell_r)
       \in(\mathcal P_X\setminus\{p,q\})^r\\
     \ell_1,\ldots,\ell_r\ \mathrm{pairwise\ distinct}}}
   |\Omega_p\cap\Omega_q\cap
     \Omega_{\ell_1}\cap\cdots\cap\Omega_{\ell_r}|. \notag
\end{align}

\begin{theorem}[The higher pair--Palm hierarchy]
\label{thm:pairpalm-hierarchy}
For every fixed $k\geq3$,
\begin{equation}\label{eq:pairpalm-anchor-identity}
 \boxed{F_k(X)=\sum_{p\ne q}T_{p,q}^{[k-2]}.}
\end{equation}
More generally, for arbitrary nonnegative numbers $b_{p,q}$, put
\[
 \mathscr P_{k,b}(X)
 =\sum_{p\ne q}
   \bigl(T_{p,q}^{[k-2]}-J_{p,q}b_{p,q}\bigr)_+.
\]
Then the exact sandwich
\begin{equation}\label{eq:pairpalm-general-sandwich}
 \mathscr P_{k,b}(X)
 \leq F_k(X)
 \leq\mathscr P_{k,b}(X)+\sum_{p\ne q}J_{p,q}b_{p,q}
\end{equation}
holds.

In particular, the periodic Poisson baseline
$b_{p,q}=\lambda_{p,q}^{k-2}$ gives
\begin{align}
 \mathscr P_k(X)
 &:=\sum_{p\ne q}
   \bigl(T_{p,q}^{[k-2]}
          -J_{p,q}\lambda_{p,q}^{k-2}\bigr)_+,
                                                        \label{eq:pairpalm-Pk}\\
 \mathscr P_k(X)
 &\leq F_k(X)
 \leq\mathscr P_k(X)+4M\lambda^k.                    \label{eq:pairpalm-periodic-sandwich}
\end{align}

There is also a finite-marginal version.  Let
\[
 A_p=|\Omega_p|,\qquad \rho_p=A_p/M,\qquad
 \mu_{p,q}=\sum_{\ell\notin\{p,q\}}\rho_\ell,
\]
and replace $\lambda_{p,q}^{k-2}$ in
\eqref{eq:pairpalm-Pk} by $\mu_{p,q}^{k-2}$, calling the result
$\mathscr P_k^{\rm fin}(X)$.  Then
\begin{equation}\label{eq:pairpalm-finite-sandwich}
 \mathscr P_k^{\rm fin}(X)
 \leq F_k(X)
 \leq\mathscr P_k^{\rm fin}(X)
   +4(1+2/X)^{k-2}M\lambda^k.
\end{equation}
Consequently $(\mathrm{HM})_k$ is equivalent to any one of
\begin{equation}\label{eq:pairpalm-equivalence}
 F_k(X)\ll X^{2+o(1)}\lambda_X^k,
 \qquad
 \mathscr P_k(X)\ll X^{2+o(1)}\lambda_X^k,
 \qquad
 \mathscr P_k^{\rm fin}(X)\ll X^{2+o(1)}\lambda_X^k.
\end{equation}
\end{theorem}

\begin{proof}
Fix a row $m$ and write $t=K(m)$.  It supplies $(t)_2$ ordered
anchor pairs $(p,q)$.  After an anchor pair is fixed, it supplies
$(t-2)_{k-2}$ ordered, pairwise distinct extension tuples.  Hence its
total contribution to the right side of
\eqref{eq:pairpalm-anchor-identity} is
\[
 (t)_2(t-2)_{k-2}=(t)_k.
\]
This proves the identity, with no missing binomial or factorial factor.

For nonnegative $A,B$ one has
$(A-B)_+\leq A\leq B+(A-B)_+$.  Apply this with
$A=T_{p,q}^{[k-2]}$ and $B=J_{p,q}b_{p,q}$, then sum over ordered
pairs.  This proves~\eqref{eq:pairpalm-general-sandwich}.
For the periodic baseline,
\[
 \sum_{p\ne q}J_{p,q}\lambda_{p,q}^{k-2}
 \leq\lambda^{k-2}\sum_{p\ne q}J_{p,q}
 =\lambda^{k-2}F_2(X)
 \leq4M\lambda^k
\]
by Lemma~\ref{lem:pairpalm-hm2-four}, proving
\eqref{eq:pairpalm-periodic-sandwich}.

For every $p$, the number of representatives below $M$ of each residue
class differs from $M/p$ by less than one.  Thus
\[
 |A_p-MZ(p)/p|\leq Z(p),
 \qquad
 \left|\sum_p\rho_p-\lambda\right|
 \leq\frac1M\sum_pZ(p)\leq\frac{2\lambda}{X}.
\]
In particular $\sum_p\rho_p\leq(1+2/X)\lambda$.
Using this in~\eqref{eq:pairpalm-general-sandwich} proves
\eqref{eq:pairpalm-finite-sandwich}.
Each positive excess is at most $F_k$, while each upper sandwich differs
from it by $O_k(M\lambda^k)$; this proves every implication in
\eqref{eq:pairpalm-equivalence}.  If $\lambda=0$, all zero sets and all
displayed quantities vanish, so the zero case requires no division.
\end{proof}

\begin{remark}[Bernoulli versus Poisson baseline]
The exact independent-Bernoulli factorial baseline is
\begin{equation}\label{eq:pairpalm-bernoulli-baseline}
 \Lambda_{p,q}^{[r]}
 =\sum_{\substack{\ell_1,\ldots,\ell_r\notin\{p,q\}\\
                   \ell_i\ \mathrm{pairwise\ distinct}}}
   \prod_{i=1}^r\frac{Z(\ell_i)}{\ell_i}
 =r!\,e_r\left(
   \left(\frac{Z(\ell)}\ell\right)_{\ell\notin\{p,q\}}
   \right).
\end{equation}
It satisfies
$0\leq\Lambda_{p,q}^{[r]}\leq\lambda_{p,q}^r$.
Therefore it may replace the Poisson baseline in
Theorem~\ref{thm:pairpalm-hierarchy}, with the same $4M\lambda^k$
error.  The Poisson power is slightly larger because its expansion
contains repeated prime labels; it is chosen in
\eqref{eq:pairpalm-Pk} because it produces the weakest positive-excess
condition among these two natural baselines.
\end{remark}

\begin{corollary}[The exact remaining pair--Palm lemma for $(\mathrm{HM})_7$]
\label{cor:pairpalm-hm7-target}
The following aggregate statement is equivalent to $(\mathrm{HM})_7$:
\begin{equation}\tag{$\mathrm{PP}_7$}\label{eq:pairpalm-PP7}
 \boxed{
 \sum_{p\ne q}
 \left(
  \sum_{m\in\Omega_p\cap\Omega_q}(K(m)-2)_5
  -J_{p,q}\lambda_{p,q}^5
 \right)_+
 \ll X^{2+o(1)}\lambda_X^7.}
\end{equation}
In particular,~\eqref{eq:pairpalm-PP7} implies
\[
 \log G_n\ll n^{20/21+o(1)}=o(n)
\]
for every $n$.  Replacing the fifth extension moment by the sixth and
$\lambda_X^7$ by $\lambda_X^8$ gives the equivalent target for
$(\mathrm{HM})_8$ and the exponent $11/12$.
\end{corollary}

\begin{proof}
The left side of~\eqref{eq:pairpalm-PP7} is exactly
$\mathscr P_7(X)$.  Theorem~\ref{thm:pairpalm-hierarchy} gives
\[
 \mathscr P_7(X)\leq F_7(X)
 \leq\mathscr P_7(X)+4X^2\lambda_X^7.
\]
This proves both directions of the equivalence.  The pointwise bounds
then follow from Theorem~\ref{thm:hm-pointwise}.
\end{proof}

The importance of Corollary~\ref{cor:pairpalm-hm7-target} is diagnostic:
it locates the missing arithmetic at the fifth factorial extension of
an actual two-prime CRT locus.  It is not itself a new cancellation
estimate and must not be advertised as weaker than $(\mathrm{HM})_7$;
the two statements are exactly equivalent.  A uniform estimate for
each individual pair would be a genuinely stronger hypothesis.

For $J_{p,q}>0$, the exact conditional content can be recorded by the
pair-Palm generating function
\begin{equation}\label{eq:pairpalm-pgf}
 \Phi_{p,q}(z)
 =\frac1{J_{p,q}}
  \sum_{m\in\Omega_p\cap\Omega_q}
  \prod_{\ell\in\mathcal P_X\setminus\{p,q\}}
       (1+zf_\ell(m)).
\end{equation}
Its fifth derivative is
\[
 \Phi_{p,q}^{(5)}(0)=\frac{T_{p,q}^{[5]}}{J_{p,q}},
\]
while the Poisson comparison $e^{\lambda_{p,q}z}$ has fifth derivative
$\lambda_{p,q}^5$ at zero.  Thus~\eqref{eq:pairpalm-PP7} says exactly
\begin{equation}\label{eq:pairpalm-pgf-target}
 \sum_{p\ne q}J_{p,q}
 \left(\Phi_{p,q}^{(5)}(0)-\lambda_{p,q}^5\right)_+
 \ll X^{2+o(1)}\lambda_X^7,
\end{equation}
with a zero summand when $J_{p,q}=0$.  Controlling only $J_{p,q}$
re-estimates the mass of the conditioning locus; it does not control this
fifth conditional coefficient.

\subsection*{The sharp deterministic envelope from present inputs}

Let
\[
 L_X=\#\{p\in\mathcal P_X:Z(p)>0\}.
\]

\begin{proposition}[Lifting the proved second moment]
\label{prop:pairpalm-unconditional-envelope}
For every fixed $k\geq3$,
\begin{align}
 F_k(X)
 &\leq(L_X-2)_{k-2}F_2(X)                              \label{eq:pairpalm-F2-lift}\\
 &\leq4X^2\lambda_X^2(L_X-2)_{k-2},                    \label{eq:pairpalm-F2-lift2}\\
 F_k(X)
 &\leq
 \sum_{\substack{p_1,\ldots,p_k\in\mathcal P_X\\
                   p_i\ \mathrm{pairwise\ distinct}}}
 \prod_{i=1}^k Z(p_i)
 \leq(2X\lambda_X)^k.                                 \label{eq:pairpalm-tuple-bound}
\end{align}
Here the first expression is interpreted as zero when $L_X<k$.
Since
\[
 L_X\ll X\min\{\lambda_X,1/\log X\},
\]
one obtains the parameter-sensitive unconditional estimate
\begin{equation}\label{eq:pairpalm-current-best}
 \boxed{
 F_k(X)\ll_k
 X^k\lambda_X^2
 \min\{\lambda_X,1/\log X\}^{k-2}.}
\end{equation}
Using the proved $Z(p)\ll p^{2/3}$ bound, hence
$\lambda_X\ll X^{2/3}/\log X$, gives
\begin{equation}\label{eq:pairpalm-current-X-bound}
 F_k(X)\ll_k\frac{X^{k+4/3}}{(\log X)^k}.
\end{equation}
In particular, the present unconditional bounds are
\[
 F_7(X)\ll\frac{X^{25/3}}{(\log X)^7},
 \qquad
 F_8(X)\ll\frac{X^{28/3}}{(\log X)^8}.
\]
\end{proposition}

\begin{proof}
Pointwise, $K(m)\leq L_X$ and
\[
 (K(m))_k=(K(m))_2(K(m)-2)_{k-2}
 \leq(K(m))_2(L_X-2)_{k-2}.
\]
Summing and using Lemma~\ref{lem:pairpalm-hm2-four} proves
\eqref{eq:pairpalm-F2-lift}--\eqref{eq:pairpalm-F2-lift2}.

Expanding $F_k$ over ordered distinct prime labels and their zero
residues, the product $p_1\cdots p_k>X^k>M$ shows that each residue
tuple has at most one representative below $M$.  This proves the first
inequality in~\eqref{eq:pairpalm-tuple-bound}; the second follows from
$\sum_pZ(p)\leq2X\lambda_X$.

Every active prime contributes at least one zero, so
$L_X\leq\sum_pZ(p)\leq2X\lambda_X$.  The prime number theorem also gives
$L_X\leq\pi(2X)-\pi(X)\ll X/\log X$.  Substitution in
\eqref{eq:pairpalm-F2-lift2} proves
\eqref{eq:pairpalm-current-best}.  Finally
Proposition~\ref{prop:zp-bound} and prime counting give
$\lambda_X\ll X^{2/3}/\log X$, yielding
\eqref{eq:pairpalm-current-X-bound}.
\end{proof}

\begin{proposition}[The anchored star is sharp for every fixed moment]
\label{prop:pairpalm-star-all-k}
The family constructed in Proposition~\ref{prop:hm3-anchored-star}
satisfies, for every fixed $k\geq3$,
\begin{equation}\label{eq:pairpalm-star-all-k}
 F_k^*(X)\asymp_k\frac{X^k}{(\log X)^k}
 \asymp_k X^k(\lambda_X^*)^k,
 \qquad
 \mathscr P_k^*(X)\asymp_k F_k^*(X).
\end{equation}
It therefore violates every estimate
\[
 F_k^*(X)\ll X^{2+\theta+o(1)}(\lambda_X^*)^k
\]
with fixed $\theta<k-2$.
\end{proposition}

\begin{proof}
In that construction $L\asymp X/\log X$ active columns all contain the
anchor row $m_0$.  Hence
\[
 F_k^*(X)\geq(K_X^*(m_0))_k=(L)_k
 \asymp_k X^k/(\log X)^k.
\]
On the other hand $\lambda_X^*\asymp1/\log X$, and the universal tuple
bound~\eqref{eq:pairpalm-tuple-bound} supplies the matching upper bound.
The sandwich~\eqref{eq:pairpalm-periodic-sandwich} also gives
\[
 \mathscr P_k^*(X)
 \geq F_k^*(X)-4X^2(\lambda_X^*)^k.
\]
For $k\geq3$ the subtracted term is lower order than the displayed
$X^k/(\log X)^k$ scale.  This proves~\eqref{eq:pairpalm-star-all-k} and
the final assertion.
\end{proof}

Thus~\eqref{eq:pairpalm-current-best} is of the correct order for the
full package of presently known primewise and pairwise incidence facts.
Any improvement to the independence scale must use an arithmetic
cross-prime restriction which the anchored-star model does not satisfy.

\subsection*{Why double Cauchy--Schwarz does not close the gap}

\begin{proposition}[Exact overlap audit]
\label{prop:pairpalm-CS-audit}
For integers $K\geq0$ and $k\geq2$,
\begin{equation}\label{eq:pairpalm-overlap-identity}
 (K)_k^2
 =(K)_2\sum_{i=0}^{k-2}\binom{k-2}{i}^2i!\,
       (K)_{2k-2-i}.
\end{equation}
Consequently Cauchy--Schwarz gives the valid inequality
\begin{equation}\label{eq:pairpalm-correct-CS}
 F_k(X)^2
 \leq F_2(X)
 \sum_{i=0}^{k-2}\binom{k-2}{i}^2i!\,
       F_{2k-2-i}(X).
\end{equation}
For $k=4$ this is
\begin{equation}\label{eq:pairpalm-correct-CS4}
 F_4^2\leq F_2(F_6+4F_5+2F_4),
\end{equation}
not $F_4^2\leq F_2F_6$.
The latter inequality and the proposed bound
$F_6\ll X^2\lambda_X^6$ from CRT injectivity are both false as
deterministic deductions.
\end{proposition}

\begin{proof}
The standard overlap identity
\[
 (Y)_a(Y)_b
 =\sum_{i=0}^{\min(a,b)}\binom ai\binom bi i!\,
   (Y)_{a+b-i}
\]
counts two ordered injective lists by the size~$i$ of their common
image.  Apply it with $Y=K-2$ and $a=b=k-2$, then multiply by
$(K)_2^2$; this is exactly
\eqref{eq:pairpalm-overlap-identity}.  Writing the right side as
$(K)_2B_K$ and summing
$(K)_k=\sqrt{(K)_2B_K}$ over rows, Cauchy--Schwarz proves
\eqref{eq:pairpalm-correct-CS}.  Formula
\eqref{eq:pairpalm-correct-CS4} is the specialization $k=4$.

For a configuration with one row of size $K=6$ and all other rows
empty,
\[
 F_2=30,\qquad F_4=360,\qquad F_6=720,
\]
so $F_4^2=6F_2F_6>F_2F_6$.  Already $K=4$ gives $F_4=24$ but
$F_6=0$.  The missing $F_5$ and $F_4$ terms in
\eqref{eq:pairpalm-correct-CS4} are precisely the overlap patterns of
the two extension tuples.

Finally, CRT injectivity for six primes says only that each ordered
choice of six prime labels and six zero residues contributes at most
one representative.  It therefore proves
\[
 F_6(X)\leq
 \sum_{p_1,\ldots,p_6\ \mathrm{distinct}}
 \prod_{i=1}^6Z(p_i)
 \leq(2X\lambda_X)^6,
\]
not $X^2\lambda_X^6$.  The missing factor $X^{-4}$ would be the
unproved assertion that the unique CRT representative lands below
$X^2$ with its expected probability.  Proposition
\ref{prop:pairpalm-star-all-k} has
$F_6^*(X)\asymp X^6(\lambda_X^*)^6$ and gives a counterexample which
also satisfies all structural conditions listed in
Proposition~\ref{prop:hm3-anchored-star}.
\end{proof}

The overlap terms cannot be suppressed by assuming a sixth-moment bound.
For an abstract sparse degree sequence, take
$R=\lfloor M\lambda^2/20\rfloor$ rows of degree~$5$ and rows of degree at
most one elsewhere.  Then
\[
 F_2=20R\leq M\lambda^2,
 \qquad F_6=0,
 \qquad F_4=120R\asymp M\lambda^2,
\]
so $(\mathrm{HM})_4$ fails by the factor $\lambda^{-2}$ as
$\lambda\to0$.  This is only an interpolation counterexample; the anchored
star supplies the stronger modular reflection-pair obstruction.

\paragraph{What the common integer sequence currently supplies.}
Lucas does give a genuine coupling:
\[
 m\bmod p\in\mathcal Z_p\quad\Longrightarrow\quad p\mid b_m.
\]
Hence the product of the distinct primes counted by $K_X(m)$ divides
$b_m$, and the known exponential growth $\log b_m=O(m)$ gives
\[
 K_X(m)\log X\leq\log b_m=O(m).
\]
For $m\asymp X^2$ this is $K_X(m)=O(X^2/\log X)$, weaker than the
trivial column bound $K_X(m)\ll X/\log X$.  Thus integer height plus
CRT does not make the joint residues uniform.  A crystalline argument
could still prove~\eqref{eq:pairpalm-PP7}, but it would need a new
quantitative equidistribution or orbit-selection theorem; growth alone
does not provide it.

Indeed the single integral sequence
$c_m=\operatorname{lcm}(1,\ldots,m)$ has $\log c_m=O(m)$, but every
prime in $(X,2X]$ divides $c_m$ as soon as $m\geq2X$.  Exponential height
and a common integer sequence can therefore coexist with maximal
cross-prime alignment.

There is also a discrepancy-scale obstruction before height enters.  For
distinct $p_1,\ldots,p_k$, CRT gives exact independence over the complete
period $P=p_1\cdots p_k$:
\begin{equation}\label{eq:pairpalm-full-period}
 \frac1P\sum_{m<P}\prod_{i=1}^kf_{p_i}(m)
 =\prod_{i=1}^k\frac{Z(p_i)}{p_i}.
\end{equation}
For $k=7$, however, $M/P<X^{-5}$.  The interval in
$(\mathrm{HM})_7$ is an extremely short prefix of one period, and
\eqref{eq:pairpalm-full-period} gives no prefix-discrepancy estimate.

Finally, primewise Fourier control alone does not eliminate the obstruction.
For the anchored star, every active zero set has two elements, so
\[
 \left|\sum_{r\in\mathcal Z_p^*}e^{2\pi iar/p}\right|\leq2
 \qquad(a\in\mathbb F_p),
\]
yet all columns meet at $m_0$ and
\eqref{eq:pairpalm-star-all-k} holds.  A large-sieve route must therefore
control mixed phases across distinct primes, not merely the Fourier
transform of each column separately.

\begin{corollary}[Threshold under an improved zero-set exponent]
\label{cor:pairpalm-alpha-threshold}
Suppose $Z(p)\ll p^{\alpha+o(1)}$ for some fixed
$0\leq\alpha<1$ and suppose
$(\mathrm{HM})_k$ holds.  Then, for every $n$,
\begin{equation}\label{eq:pairpalm-alpha-threshold}
 \log G_n\ll n^{\max\{2/3,\,\alpha+2/k\}+o(1)}.
\end{equation}
In particular, a bound $Z(p)\ll p^{1/2+o(1)}$ makes
$(\mathrm{HM})_5$ sufficient for the full conjecture.
\end{corollary}

\begin{proof}
Prime counting gives $\lambda_X\ll X^\alpha/\log X$.
The maximum-row argument in the proof of
Theorem~\ref{thm:hm-pointwise} then yields
\[
 \max_{m<X^2}K_X(m)
 \ll X^{2/k+o(1)}\lambda_X+1
 \ll X^{\alpha+2/k+o(1)}/\log X.
\]
The same dyadic summation gives the exponent $\alpha+2/k$ for the
lower-digit channel.  Combining it with the already proved small-prime and
$O(n^{2/3})$ companion-block bounds proves
\eqref{eq:pairpalm-alpha-threshold}.
For $\alpha=1/2$ and $k=5$ the exponent is $9/10<1$.
\end{proof}

\subsection*{Finite verification and stall report}

The standard-library script \texttt{pairpalm\_verify.py} performs seven
groups of exact checks.  It computes the Ap\'ery zero sets both from the
divided recurrence and from the cleared recurrence $y_m=(m!)^3b_m$, and
cross-checks fixed zero sets, reflection, and no consecutive zeros.  For
the actual block $X=128$ it finds
\[
 \#\mathcal P_X=23,\qquad
 \#\{p:Z(p)>0\}=13,\qquad
 \sum_pZ(p)=30,
\]
\[
 \#\{m:K(m)=0,1,2,3\}=(13976,2242,156,10),
 \qquad
 (F_2,F_3,F_4,\ldots,F_8)=(372,60,0,\ldots,0).
\]
At the first tested block with a fourth-order collision, $X=1024$, two
independent algorithms---row scattering and direct CRT enumeration in
prime/residue space---agree on
\[
 \#\mathcal P_X=137,\qquad
 \#\{p:Z(p)>0\}=47,\qquad
 \sum_pZ(p)=120,\qquad \lambda_X=0.082244154\ldots,
\]
\[
 (F_2,F_3,F_4,F_5,F_6,F_7,F_8)
 =(6894,486,24,0,0,0,0).
\]
The existential-third-prime aggregate is $474$, rather than the correct
multiplicity count $F_3=486$.  The periodic fourth-moment Palm excess is
$23.448043852\ldots$ (the finite-marginal value is
$23.448002897\ldots$), and every periodic and finite sandwich is checked
with exact rational arithmetic for $3\leq k\leq8$.

A separate synthetic incidence table has a row of size~$8$ and gives
nonzero strict positive-excess tests through $k=8$.  The script also checks
the overlap identity~\eqref{eq:pairpalm-overlap-identity}, both
Cauchy--Schwarz counterexamples, the CRT six-tuple failure, the height
counterexample, and the exact moments through $k=8$ of the anchored star at
$X=64$:
\[
 F_7^*=3{,}991{,}680,
 \qquad F_8^*=19{,}958{,}400.
\]
All seven groups print \texttt{PASS}.  The vanishing of the actual
$F_7,F_8$ at these finite scales tests identities only; it is not evidence
for the required uniform asymptotic bound.

No estimate of~\eqref{eq:pairpalm-PP7}, $(\mathrm{HM})_7$, or
$(\mathrm{HM})_8$ is proved here.  The sharp unconditional output from
the present toolbox is Proposition
\ref{prop:pairpalm-unconditional-envelope}; the anchored-star family
shows why the known reflection, interval, fiber, energy, gap-certificate,
and row-codegree statements cannot improve it.  The exact remaining
load-bearing input is the aggregate fifth-extension positive-excess
bound~\eqref{eq:pairpalm-PP7}, or an arithmetic theorem which implies it
such as a suitable cross-prime Fourier/dispersion estimate.
